\section{Experimental Setup}
\label{sec:experiments}

\subsection{Feature Set}

The full feature vector comprises 22 dimensions organized into three
families:
\begin{itemize}
    \item \textbf{LLM features (7d):} BGE-large cosine similarity,
          BM25 score, Jaccard coefficient, character n-gram overlap,
          and neighborhood-derived similarity measures.
    \item \textbf{Structural features (13d):} Shortest-path distance,
          common-neighbor count, neighbor-set Jaccard index, Adamic-Adar
          index, bridge centrality, degree-product features, and
          framework-pair indicators derived from the four-source
          knowledge graph.
    \item \textbf{Opus features (2d):} Claude Opus tier prediction
          (0--3) and associated confidence score.
\end{itemize}

Figure~\ref{fig:feature_distributions} shows the distribution of six
representative features stratified by expert tier, illustrating the
discriminative signal available in each family.
Figure~\ref{fig:correlation_matrix} presents the Pearson correlation
matrix across all 22 features, revealing that the three families
capture largely orthogonal signals.

\begin{figure}[t]
\centering
\includegraphics[width=\linewidth]{fig_feature_distributions}
\caption{Distribution of six representative features stratified by
expert tier label. Each panel shows violin and strip plots for one
feature across the four tiers, illustrating the per-family
discriminative signal.}
\label{fig:feature_distributions}
\end{figure}

\begin{figure}[t]
\centering
\includegraphics[width=\linewidth]{fig_correlation_matrix}
\caption{Pearson correlation matrix of all 22 v6 features on the test
set ($n=400$). Black lines delineate the three feature families
(LLM, Structural, Opus). Cross-family correlations are generally low,
confirming that the three signal families capture complementary
information.}
\label{fig:correlation_matrix}
\end{figure}

\subsection{Training Configuration}

\begin{itemize}
    \item \textbf{Classifier:} LightGBM multiclass (softmax objective).
    \item \textbf{Hyperparameters:} \texttt{n\_estimators}${}=500$,
          learning rate${}=0.05$,
          \texttt{min\_samples\_leaf}${}=5$.
    \item \textbf{Class weights:} Inverse-frequency weighting to
          compensate for class imbalance.
    \item \textbf{Training data:} 1,515 LLM-SME labeled pairs.
    \item \textbf{Calibration data:} 150 held-out pairs for Mondrian
          conformal calibration.
    \item \textbf{Test data:} 400 expert-labeled frozen pairs
          (\texttt{v1\_frozen}), never seen during training or tuning.
          SHA-256 hash verification ensures immutability.
\end{itemize}

\subsection{Evaluation Metrics}

We report four accuracy metrics at increasing granularity:
\begin{itemize}
    \item \textbf{Tier accuracy:} Exact-match accuracy across the four
          tiers.
    \item \textbf{Adjacent accuracy:} Fraction of predictions within
          $\pm 1$ tier of the ground-truth label, capturing near-miss
          performance on an ordinal scale.
    \item \textbf{Binary accuracy:} Accuracy after collapsing tiers
          into a binary ``related or stronger'' vs.\ ``partial or
          weaker'' distinction.
    \item \textbf{Macro F1:} Unweighted average of per-class F1 scores,
          treating all four tiers equally regardless of prevalence.
\end{itemize}

Additionally, we report per-class precision, recall, and F1, as well
as Mondrian conformal coverage and mean prediction set size at
$\alpha = 0.10$. All confidence intervals are computed via
10,000-iteration bootstrap resampling.
